\chapter{La problématique de financement de l'open source}

Beacoup de projets libres sont gratuits, mais cela ne signifie pas forcément que
rien n'a été investi : il existe de nombreux modèles de financement des projets
opensource dont certains initiés par des sociétés comme RedHat ou Apple qui sont
en croissance constante.\\
Le logiciel libre peut donc être rentable à condition de savoir exploiter
le modèle de financement adapté à l'ampleur du projet.

\section{Quelques modèles alternatifs}

    \subsection{Les dons}

Le financement de projets par dons est très répandu.
Si ce modèle peut servir d'appoint aux développeurs, il n'est viable que dans
très peu de cas.
    
    \subsection{donations par micropaiement}

Le modèle du pay per like est une version plus organisée du modèle par dons. Les
dons se font non plus directement au projet mais passent par un oganisme chargé
de répartir équitablement un budget mensuel, défini par le donneur, entre tous
les projets qu'il aura explicitement marqués dans le mois.
\url{flattr.com}, par exemple, permet distruibuer le budget entre les projets
marqués sur \url{github.com}.\\
Flattr, bien que compatible pour le financement de projets informatiques est
surtout prévu pour récompenser les créateurs de contenu artistique.

    \subsection{Licence Double}

La licence double est un modèle adopté notamment par MySQL consiste à permettre
l'utilisation de son produit différemment suivant l'usage. Dans le cas de MySQL,
l'utilisation personnelle ou l'inclusion dans un projet sous licence GPL est
gratuite tandis qu'une utilisation profesionnelle est payante.

    \subsection{Modèle Redhat}

Redhat distribue son produit gratuitement mais propose un support payant
garantissant la stabilité de son produit. Ce modèle a été adopté par
l'entreprise en 2003, qui depuis cette date est en constante croissance.
C'est un modèle similaire qui régit Canonical.

    \subsection{Modèle Apple}

Le système d'exploiation OSX d'apple est dérivé d'un système ouvert.\\
La stratégie d'apple pour son système d'exploitation OSX a consisté à changer la
licence et apporter des modifications pour ensuite le vendre sur ses machines.

    \subsection{Modèle Bountysource}

\url{bountysource.com} est un site d'annonces de demande de code. Les
utilisateurs peuvent demander des fonctionnalités et se regrouper à plusieurs
sur une même annonce pour augmenter la prime. Ainsi, plus une fonction intéresse
d'utilisateurs, plus la somme proposée pour son implémentation est élevée et
donc susceptible d'intéresser un développeur.

\section{Gittip confronté aux autres modèles}

Gittip est un modèle proche de celui de Flattr à la différence que les
donnations sont obligatoirement récurrentes et faites à un développeur ou une
équipe plutot qu'un projet précis.

\subsection{Les modèles proches}

\paragraph{}
Le caracère récurrent des dons assure plus de stabilité que les dons ponctuels
via Flattr ou Paypal.

\paragraph{}
Les équipes financées grace a Gittip dépendent d'un grand nombre d'acteurs,
ce qui leur permet de mesurer leur popularité et de ne pas dépendre d'un seul
investisseur qui pourrait décider d'abandonner le projet.

\paragraph{}
Gittip est lui même financé via Gittip et non via le prélèvement sur les dons
fait aux développeurs, le seuls prélèvements effectués servent à couvrir
les opérations bancaires.

\begin{figure}[h!]
\caption{Comparatif des prélèvements pas plateforme}
\begin{center}
\begin{tabular}{|c|c|c|}
    \hline % Permet d'avoir une séparation via un trait
    \textbf{Paypal} & \textbf{Gittip} & \textbf{Flattr} \\
    \hline
    2,2\% + 30\textcent & 2,9\% + 30\textcent & 10\% \\
    \hline
\end{tabular}
\end{center}
\end{figure}

Dans le tableau ci-dessus, le prix des transactions pour utilisé Paypal est
celui appliqué au bouton "Donate".

\paragraph{}
Gittip permet également de gérer efficacement son budget à l'interieur d'une
équipe grâce au mécanisme précédemment décrit, ce que ne gèrent pas flattr ni
Paypal.

\subsection{Les autres modèles}

Les modèles restants visent soit à lancer un projet, soit à financer des projets
de grande ampleur, et donc nécessitant un investissement plus important.
Actuellement, Chad Whitacre gagne 23300\$ par an alors que le salaire moyen de
développeur est de 90000\$ par an. Dans l'état actuel des choses, Gittip n'est
donc pas viable en tant que source principale de revenus. Toutefois, il est
important de noter que la somme échangée par semaine augmente avec le nombre
d'utilisateurs actifs depuis la création du site.

\begin{figure}[h!]
  \caption{Nombre d'utilisateurs en fonction du temps}
  \centering
    \includegraphics[width=16cm]{images/gittip_amount.eps}
\end{figure}
\begin{figure}[h!]
  \caption{Somme échangée au sein de Gittip (en \$) en fonction du temps}
  \centering
    \includegraphics[width=16cm]{images/gittip_user.eps}
\end{figure}

parler de : pas de financement important (comparé a kickstarter pour initier un projet ou
 le financement par une société tierce comme chez mozilla)


